\begin{landscape}
\pdfbookmark[2]{Zweite Rechnung, [2]}{pallas-IV-fragment2}
\begin{center}[2.]\end{center}
\begingroup
\setlength{\lineskip}{4pt}\setlength{\lineskiplimit}{4pt}
\setlength{\arraycolsep}{3pt}
\newcommand{\IVtwoDayComma}{\mathord{\rlap{\(^{\mathrm d}\)}{,}\kern.2em}}
\newcommand{\IVtwoLowNode}{\mathord{\underset{\smash{\raisebox{.4ex}{\rule{.28em}{.3pt}}}}{\AscendingNode}}}
\newcommand{\IVtwoBeforeVon}{\makebox[0pt][r]{\raisebox{.05ex}{\tiny .}\kern.12em}}
\noindent\hspace*{\parindent}[Für die einzelnen Oppositionen ergeben sich aus unserer Tafel die folgenden Werte für die Störungen:
\begin{center}
\begingroup\setlength{\tabcolsep}{7pt}
\begin{tabular}{@{}l r|r|r|r|r|r|r|r@{}}
\multicolumn{2}{c|}{\rule[-2pt]{0pt}{10pt}Opposition} & \(\delta i\) & \(\delta\AscendingNode\) & \(50\delta n\) & \(\delta\!\int n\mathrm{d}t\) & \(\delta\varphi\) & \(\delta\tilde\omega\) & \(\delta\varepsilon\)\\
\hline
I. & \(181\IVtwoDayComma02\) & \(-0'',05\) & \(+0'',45\) & \(-0'',601\) & \(-0'',10\) & \(-3'',61\) & \(+0'',02\) & \(+4'',68\)\\
II. & \(608,21\) & \(-11,17\) & \(+17,88\) & \(-17,204\) & \(-71,39\) & \(-77,82\) & \(-75,37\) & \(+39,23\)\\
III. & \(1064,47\) & \(-34,58\) & \(+0,81\) & \(-29,977\) & \(-300,94\) & \(-112,39\) & \(-275,95\) & \(-4,82\)\\
IV. & \(1585,61\) & \(-21,83\) & \(-12,48\) & \(-11,825\) & \(-529,81\) & \(-160,90\) & \(-420,12\) & \(+0,99\)\\
V. & \(2034,89\) & \(-16,26\) & \(-43,55\) & \(-28,115\) & \(-659,68\) & \(-280,59\) & \(-499,69\) & \(+147,52\)\\
VI. & \(2457,68\) & \(-67,39\) & \(-46,12\) & \(-123,425\) & \(-1271,15\) & \(-674,35\) & \(-1325,85\) & \(+10,77\)\\
VII. & \(2974,81\) & \(-152,12\) & \(-163,89\) & \(-94,923\) & \(-2421,91\) & \(-834,09\) & \(-1163,53\) & \(-179,74\)
\end{tabular}\endgroup
\end{center}
\par
Wir nehmen nun als genäherte Werthe für die Normalelemente, welche wieder etwa für \(1803.156^{\mathrm d}\) gelten
\[\begin{aligned}
i&=34^0\,37'\,46'',0&\qquad\varphi&=14^0\,12'\,47'',5\\
\AscendingNode&=172\quad28\quad29,5&\tilde\omega&=121\quad8\quad0
\end{aligned}\]
\[n=770'',7372\]
und berechnen zur Bestimmung ihrer Correctionen aus den beobachteten heliocentrischen Längen, gerade \hypertarget{native-493-cont}{}\label{native-493-cont}wie in der ersten Rechnung, die Werthe der Epoche der mittlern Länge wie folgt:
\begingroup\small
\[\begin{array}{@{}l r|r|r|r|r|r|r@{}}
\alpha&277^0\,38'\,59'',17&336^0\,59'\,12'',67&67^0\,18'\,16'',89&223^0\,33'\,50'',20&303^0\,58'\,20'',58&359^0\,34'\,27'',29&152^0\,41'\,27'',75\\
\AscendingNode&172\;28\;29,95&172\;28\;47,38&172\;28\;30,31&172\;28\;17,02&172\;27\;45,95&172\;27\;43,38&172\;25\;45,61\\
i&34\;37\;45,95&34\;37\;34,83&34\;37\;11,42&34\;37\;24,17&34\;37\;29,74&34\;36\;38,61&34\;35\;13,88\\
\delta&+33\;41\;4,98&+10\;27\;3,72&-33\;40\;32,87&+28\;14\;54,73&+27\;20\;32,29&-4\;53\;1,98&-13\;6\;31,65\\
\hline
v&275\;3\;20,53&333\;51\;47,85&69\;53\;47,97&228\;52\;48,72&298\;31\;39,95&1\;4\;57,17&148\;52\;48,46\\
\tilde\omega&121\;8\;0,02&121\;6\;44,63&121\;3\;24,05&121\;0\;59,88&129\;59\;40,31&120\;45\;54,15&120\;48\;36,47\\
\varphi&14\;12\;43,89&14\;11\;29,68&14\;10\;55,11&14\;10\;6,60&14\;8\;6,91&14\;1\;33,15&13\;58\;53,41\\
v-\tilde\omega&153\;55\;20,51&212\;45\;3,22&308\;50\;23,92&107\;51\;48,84&177\;31\;59,64&240\;19\;3,02&28\;4\;11,99\\
M&139\;10\;10,21&230\;38\;5,89&328\;18\;2,61&79\;50\;56,99&176\;3\;46,23&266\;36\;12,28&16\;53\;56,10\\
L&260\;18\;10,23&351\;44\;50,52&89\;21\;26,66&200\;51\;56,87&297\;3\;26,54&27\;22\;6,43&137\;42\;32,57\\
L_0&260\;18\;5,65&351\;45\;22,78&89\;26\;32,42&201\;0\;45,69&297\;11\;58,70&27\;43\;6,81&138\;25\;54,22\\
(\varepsilon_0)&254\;56\;43,42&254\;56\;46,66&254\;57\;15,53&254\;57\;25,66&254\;57\;39,08&254\;58\;6,90&254\;58\;17,83
\end{array}\]
\endgroup
\par
Ferner ergeben sich die Werthe
\[\begin{array}{@{}r r|r|r|r|r|r|r@{}}
m+1:&-0,499&-0,446&+0,315&-0,065&-0,595&-0,179&+0,379\\
n:&-1,210&+1,447&+1,185&-2,013&-0,124&+1,985&-0,666.
\end{array}\]
\par
Es folgen hieraus die Correctionen der Elemente:
\[\begin{aligned}
\Delta\tilde\omega&=-5'',8\\
\Delta\varphi&=-1,2\\
\Delta n&=-0,0012
\end{aligned}\]
und nach ihrer Anbringung die verbesserten Werthe
\[\begin{aligned}
\tilde\omega&=121^0\,7'\,54'',2&\qquad n&=770'',7360\\
\varphi&=14\quad12\quad46,3&\varepsilon&=254^0\,56'\,38'',19.
\end{aligned}\]
\par
Für die einzelnen Oppositionen erhält man für die Epoche:
\[\begin{array}{@{}r|r|r|r|r|r|r@{}}
254^0\,56'\,46'',86&31'',23&39'',58&36'',99&35'',04&42'',70&34'',95
\end{array}\]
\SourceLeaf{494}\noindent und es bleiben folgende Unterschiede in den mittleren Längen übrig (vgl. Bd. VI S. 329):
\[\begin{array}{@{}r|r|r|r|r|r|r@{}}
+8'',67&-6'',96&+1'',39&-1'',20&-3'',15&+4'',51&-3'',24.
\end{array}\]
\par
Zur Bestimmung der Correctionen \(\Delta i\) und \(\Delta\AscendingNode\) stellt sich die Rechnung, wie folgt:
\begingroup\small
\[\begin{array}{@{}r r|r|r|r|r|r|r@{}}
n={}&770'',7240&770'',3919&770'',1365&770'',4995&770'',1737&768'',2675&768'',8375\\
p={}&0,415\,0716&0,415\,2758&0,415\,4084&0,415\,3236&0,415\,5730&0,416\,7063&0,416\,6593\\
r={}&0,523\,2667&0,515\,5596&0,353\,3334&0,449\,2217&0,537\,0429&0,472\,2310&0,332\,7308\\
\beta-\gamma={}&12^0\,45'\,32'',66&4^0\,34'\,42'',04&20^0\,50'\,26'',13&13^0\,56'\,27'',09&10^0\,23'\,38'',28&2^0\,29'\,4'',62&10^0\,42'\,5'',34\\
\gamma={}&+33\;41\;3,34&+10\;27\;7,76&-33\;40\;28,77&+28\;14\;58,51&+27\;20\;15,42&-4\;53\;5,48&-13\;6\;13,86\\
\gamma-\delta={}&+1'',64&-4'',04&-4'',10&-3'',78&-16'',87&+3'',50&-18'',14
\end{array}\]
\endgroup
\par
Hieraus ergeben sich wieder sieben Gleichungen \IVtwoBeforeVon von der Form
\[c\Delta\AscendingNode+f\Delta i=\gamma-\delta,\]
welche in unserm Falle lauten:
\[\begin{aligned}
+0,9870\,\Delta i+0,1252\,\Delta\AscendingNode&=+1'',64\\
+0,3815\,\Delta i+0,6436\,\Delta\AscendingNode&=-4,04\\
-0,9869\,\Delta i+0,1251\,\Delta\AscendingNode&=-4,10\\
+0,8918\,\Delta i-0,3365\,\Delta\AscendingNode&=-3,78\\
+0,8725\,\Delta i+0,3610\,\Delta\IVtwoLowNode&=-16,87\\
-0,1846\,\Delta i+0,6914\,\Delta\AscendingNode&=+3,50\\
-0,4725\,\Delta i-0,6154\,\Delta\AscendingNode&=-18,14.
\end{aligned}\]
\par
Da bei der Ausführung dieser Rechnung die Beobachtungen der siebenten Opposition noch nicht vorlagen und die Breite der fünften Opposition schlecht beobachtet war, so geben wir das Resultat nur mit Berücksichtigung der vier ersten und der sechsten Gleichung, nemlich
\[\Delta i=-1'',7\qquad\Delta\AscendingNode=-8'',6,\]
wonach die corrigirten Werthe
\[i=34^0\,37'\,44'',3\qquad\AscendingNode=172^0\,28'\,20'',9\]
werden. Nach Anbringung der gefundenen Correctionen bleiben folgende Unterschiede in den geocentrischen Breiten übrig (vgl. Band VI Seite 329):
\[\begin{array}{@{}r|r|r|r|r|r|r@{}}
-1'',08&-10'',20&-3'',52&-2'',39&(-21'',42)&-2'',11&-12'',08.]
\end{array}\]
\begin{center}\rule{.22\linewidth}{.3pt}\end{center}
\endgroup
\end{landscape}
